\section{\texorpdfstring{Reading the Graph from \(f'\) and \(f''\)}{Reading the Graph from f' and f''}}

\subsection{First-Order Shape}

\begin{toolkitbox}{First-Order Shape --- Quick Reference}
\small
\begin{tabular}{l l}
\toprule
\textbf{Statement} & \textbf{Role} \\
\midrule
\hyperref[thm:zero-derivative-implies-constant]{Zero derivative} & Vanishing derivative forces constant behavior \\
\hyperref[cor:equal-derivatives-constant-difference]{Equal derivatives} & Functions differ by a constant \\
\hyperref[thm:nondecreasing-iff-nonneg-derivative]{Nonnegative derivative} & Detects nondecreasing behavior \\
\hyperref[thm:nonincreasing-iff-nonpos-derivative]{Nonpositive derivative} & Detects nonincreasing behavior \\
\hyperref[lem:positive-derivative-implies-locally-increasing]{Positive derivative} & Local increasing behavior \\
\hyperref[lem:negative-derivative-implies-locally-decreasing]{Negative derivative} & Local decreasing behavior \\
\hyperref[thm:first-derivative-test-maximum]{First derivative test: max} & Sign change detects a local maximum \\
\hyperref[thm:first-derivative-test-minimum]{First derivative test: min} & Sign change detects a local minimum \\
\bottomrule
\end{tabular}
\end{toolkitbox}

% =========================================================
% Section: Understanding Derivatives --- Monotonicity and Extrema
% =========================================================
% ---------------------------------------------------------

\begin{remark*}[Interpretation]
The Mean Value Theorem is now available, so the harvest can begin. Vanishing
derivatives force constancy, derivative signs force monotonicity, and second
derivatives detect bending when enough regularity is present. Darboux and the
smoothness hierarchy show that derivatives carry hidden rigidity even when
they are not continuous.
\end{remark*}

\begin{definitionbox}{Definition (Increasing at a Point)}
\begin{definition}[Increasing at a Point]
\label{def:increasing-at-a-point}
Let $A \subseteq \mathbb{R}$, let $f : A \to \mathbb{R}$, and let $c \in A$. We say that $f$ is \emph{increasing at $c$} if there exists $\delta > 0$ such that for all
\[
x \in A \cap (c-\delta,c)
\quad \text{and} \quad
y \in A \cap (c,c+\delta),
\]
we have $f(x) \leq f(c) \leq f(y)$.
\end{definition}
\end{definitionbox}

\begin{remark*}[Standard quantified statement]
\[
\exists \delta > 0\;\forall x \in A \cap (c-\delta,c)\;\forall y \in A \cap (c,c+\delta)
\;\bigl(f(x) \leq f(c) \leq f(y)\bigr).
\]
\end{remark*}

\begin{remark*}[Predicate reading]
\[
\exists \delta>0\;\forall x<c<y\;
\bigl(x,y\in A\cap(c-\delta,c+\delta)\Rightarrow f(x)\le f(c)\le f(y)\bigr).
\]
\end{remark*}

\begin{remark*}[Interpretation]
$f$ is increasing at $c$ if $c$ lies between the values of $f$ on either side in some punctured neighbourhood: $f$ is smaller to the left and larger to the right (weakly). This is a strictly local notion --- it says nothing about $f$ outside the $\delta$-neighbourhood of $c$. A function can be increasing at $c$ without being increasing (globally or on any interval). The non-strict inequalities permit the function to be flat on one or both sides of $c$, distinguishing this from the notion of a strict local increase.
\end{remark*}

\begin{dependencies}
\begin{itemize}
  \item \hyperref[def:function-strictly-increasing]{Definition: Strictly Increasing Function (functions chapter)}
\end{itemize}
\end{dependencies}

% ---------------------------------------------------------

\begin{definitionbox}{Definition (Decreasing at a Point)}
\begin{definition}[Decreasing at a Point]
\label{def:decreasing-at-a-point}
Let $A \subseteq \mathbb{R}$, let $f : A \to \mathbb{R}$, and let $c \in A$. We say that $f$ is \emph{decreasing at $c$} if there exists $\delta > 0$ such that for all
\[
x \in A \cap (c-\delta,c)
\quad \text{and} \quad
y \in A \cap (c,c+\delta),
\]
we have $f(x) \geq f(c) \geq f(y)$.
\end{definition}
\end{definitionbox}

\begin{remark*}[Standard quantified statement]
\[
\exists \delta > 0\;\forall x \in A \cap (c-\delta,c)\;\forall y \in A \cap (c,c+\delta)
\;\bigl(f(x) \geq f(c) \geq f(y)\bigr).
\]
\end{remark*}

\begin{remark*}[Predicate reading]
\[
\exists \delta>0\;\forall x<c<y\;
\bigl(x,y\in A\cap(c-\delta,c+\delta)\Rightarrow f(x)\ge f(c)\ge f(y)\bigr).
\]
\end{remark*}

\begin{remark*}[Interpretation]
The function is decreasing at $c$ when nearby values to the left sit above $f(c)$ and nearby values to the right sit below it.
\end{remark*}

\begin{dependencies}
\begin{itemize}
  \item \hyperref[def:increasing-at-a-point]{Definition: Increasing at a Point}
\end{itemize}
\end{dependencies}

% ---------------------------------------------------------

\begin{theorembox}{Theorem (Zero Derivative Implies Constant)}
\begin{theorem}[Zero Derivative Implies Constant]
\label{thm:zero-derivative-implies-constant}
Let $I := [a,b] \subseteq \mathbb{R}$ be a closed interval, and let $f : I \to \mathbb{R}$ be continuous on $I$ and differentiable on $(a,b)$. If
\[
f'(x) = 0 \qquad \text{for all } x \in (a,b),
\]
then $f$ is constant on $I$.
\hyperref[prf:zero-derivative-implies-constant]{Go to proof.}
\end{theorem}
\end{theorembox}

\begin{remark*}[Standard quantified statement]
\[
\forall x \in (a,b)\;\operatorname{Derivative}(0,f,x,\mathbb{R},\mathbb{R})
\;\Rightarrow\;
\exists C \in \mathbb{R}\;\forall x \in I\;(f(x) = C).
\]
\end{remark*}

\begin{remark*}[Predicate reading]
If the derivative of $f$ is zero at every interior point of the interval,
then $f$ is constant on the whole closed interval.
\[
  f\colon A\to\mathbb{R},\\
  \operatorname{Derivative}(D,f,a,\mathbb{R},\mathbb{R}).
\]
\end{remark*}

\begin{remark*}[Interpretation]
\[
\lnot\bigl(\exists C \in \mathbb{R}\;\forall x \in I\;(f(x)=C)\bigr)
\;\Rightarrow\;
\exists x \in (a,b)\;\lnot\operatorname{Derivative}(0,f,x,\mathbb{R},\mathbb{R}).
\]
If $f$ is not constant on $I$, then $f'$ is not identically zero on $(a,b)$.


\[
\lnot\operatorname{ConstantOn}(f,I) \;\Rightarrow\; \exists x \in (a,b)\;(f'(x) \neq 0).
\]


The MVT is the engine here: if $f'(x) = 0$ everywhere on $(a,b)$, then for any two points $x_1 < x_2$ in $I$, the MVT supplies a $c$ with $f(x_2)-f(x_1) = f'(c)(x_2-x_1) = 0$, so $f(x_1) = f(x_2)$. Since this holds for all pairs, $f$ is constant. The theorem is the foundation of antiderivative uniqueness: two antiderivatives of the same function differ by a constant. The contrapositive is used in ODE uniqueness arguments --- if a solution is not identically the zero solution, its derivative is not identically zero somewhere.
\end{remark*}

\begin{dependencies}
\begin{itemize}
  \item \hyperref[thm:mean-value-theorem]{Theorem: Mean Value Theorem}
\end{itemize}
\end{dependencies}

% ---------------------------------------------------------

\begin{corollarybox}{Corollary}
\begin{corollary}[Equal Derivatives Imply Difference Is Constant]
\label{cor:equal-derivatives-constant-difference}
Let $I := [a,b] \subseteq \mathbb{R}$, and let $f, g : I \to \mathbb{R}$ be continuous on $I$ and differentiable on $(a,b)$. If
\[
f'(x) = g'(x) \qquad \text{for all } x \in (a,b),
\]
then there exists $C \in \mathbb{R}$ such that $f(x) = g(x) + C$ for all $x \in I$.
\hyperref[prf:equal-derivatives-constant-difference]{Go to proof.}
\end{corollary}
\end{corollarybox}

\begin{remark*}[Standard quantified statement]
\[
\forall x \in (a,b)\;(f'(x)=g'(x))
\;\Rightarrow\;
\exists C \in \mathbb{R}\;\forall x \in I\;(f(x)=g(x)+C).
\]
\end{remark*}

\begin{remark*}[Predicate reading]
If two differentiable functions have equal derivatives at every interior
point of the interval, then their difference is constant on the interval.
\[
  f\colon A\to\mathbb{R},\\
  \operatorname{Derivative}(D,f,a,\mathbb{R},\mathbb{R}).
\]
\end{remark*}

\begin{remark*}[Interpretation]
\[
\forall C \in \mathbb{R}\;\exists x \in I\;(f(x) \neq g(x)+C)
\;\Rightarrow\;
\exists x \in (a,b)\;(f'(x) \neq g'(x)).
\]
If $f - g$ is not a constant function, then $f'$ and $g'$ disagree somewhere.


\[
\lnot\operatorname{ConstantOn}(f-g,\,I) \;\Rightarrow\; \exists x\in(a,b)\;(f'(x)\neq g'(x)).
\]


This is the corollary that makes antiderivatives well-defined up to a constant: any two antiderivatives of the same function on an interval differ by a constant. The proof is one line --- apply the previous theorem to $h = f - g$, whose derivative is $(f-g)' = f'-g' = 0$.
\end{remark*}

\begin{dependencies}
\begin{itemize}
  \item \hyperref[thm:zero-derivative-implies-constant]{Theorem: Zero Derivative Implies Constant}
\end{itemize}
\end{dependencies}

% ---------------------------------------------------------

\begin{theorembox}{Theorem (Nondecreasing Iff Nonnegative Derivative)}
\begin{theorem}[Nondecreasing Iff Nonnegative Derivative]
\label{thm:nondecreasing-iff-nonneg-derivative}
Let $I \subseteq \mathbb{R}$ be an interval, and let $f : I \to \mathbb{R}$ be differentiable on $I$. Then $f$ is nondecreasing on $I$ if and only if
\[
f'(x) \geq 0 \qquad \text{for all } x \in I.
\]
\hyperref[prf:nondecreasing-iff-nonneg-derivative]{Go to proof.}
\end{theorem}
\end{theorembox}

\begin{remark*}[Standard quantified statement]
\[
\operatorname{OrderPreserving}(f,I)
\;\iff\;
\forall x \in I\;(f'(x) \geq 0).
\]
\end{remark*}

\begin{remark*}[Interpretation]
\[
\exists x \in I\;(f'(x) < 0)
\;\Rightarrow\;
\lnot\operatorname{OrderPreserving}(f,I).
\]
If the derivative is negative at some point, the function is not nondecreasing.


\[
\exists x \in I\;(f'(x) < 0) \;\Rightarrow\; \lnot\operatorname{OrderPreserving}(f,I).
\]


The derivative is the instantaneous rate of change; requiring it to be nonnegative everywhere exactly characterises functions that never decrease. The forward direction ($f$ nondecreasing $\Rightarrow$ $f' \geq 0$) follows from the definition: the difference quotient is nonnegative when $f$ is nondecreasing, so the limit is nonnegative. The reverse direction ($f' \geq 0 \Rightarrow$ nondecreasing) is a consequence of the MVT: for $x < y$, the MVT gives $f(y)-f(x) = f'(c)(y-x) \geq 0$.

Note carefully: $f' \geq 0$ characterises \emph{nondecreasing} (weak), not \emph{strictly increasing} (strict). Strict monotonicity requires $f' > 0$ at all but isolated points. The function $f(x) = x^3$ satisfies $f'(0) = 0$ but is strictly increasing; $f' \geq 0$ with equality only at isolated points is sufficient for strict increase, but $f' > 0$ everywhere is not necessary.
\end{remark*}

\begin{dependencies}
\begin{itemize}
  \item \hyperref[def:derivative-at-a-point]{Definition: Derivative at a Point}
  \item \hyperref[thm:mean-value-theorem]{Theorem: Mean Value Theorem}
\end{itemize}
\end{dependencies}

% ---------------------------------------------------------

\begin{theorembox}{Theorem}
\begin{theorem}[Nonincreasing Iff Nonpositive Derivative]
\label{thm:nonincreasing-iff-nonpos-derivative}
Let $I \subseteq \mathbb{R}$ be an interval, and let $f : I \to \mathbb{R}$ be differentiable on $I$. Then $f$ is nonincreasing on $I$ if and only if
\[
f'(x) \leq 0 \qquad \text{for all } x \in I.
\]
\hyperref[prf:nonincreasing-iff-nonpos-derivative]{Go to proof.}
\end{theorem}
\end{theorembox}

\begin{remark*}[Standard quantified statement]
\[
\operatorname{OrderReversing}(f,I)
\;\iff\;
\forall x \in I\;(f'(x) \leq 0).
\]
\end{remark*}

\begin{remark*}[Interpretation]
\[
\exists x \in I\;(f'(x) > 0) \;\Rightarrow\; \lnot\operatorname{OrderReversing}(f,I).
\]


\[
\exists x \in I\;(f'(x) > 0) \;\Rightarrow\; \lnot\operatorname{OrderReversing}(f,I).
\]


Nonpositive derivative everywhere is exactly the infinitesimal signature of a function that never increases.
\end{remark*}

\begin{dependencies}
\begin{itemize}
  \item \hyperref[thm:nondecreasing-iff-nonneg-derivative]{Theorem: Nondecreasing Iff Nonnegative Derivative}
\end{itemize}
\end{dependencies}

% ---------------------------------------------------------

\begin{lemmabox}{Lemma}
\begin{lemma}[Positive Derivative Implies Locally Increasing]
\label{lem:positive-derivative-implies-locally-increasing}
Let $I \subseteq \mathbb{R}$ be an interval, let $f : I \to \mathbb{R}$, let $c \in I$, and suppose that $f$ is differentiable at $c$. If $f'(c) > 0$, then there exists $\delta > 0$ such that
\[
f(x) < f(c) \quad \text{for all } x \in I \text{ with } c - \delta < x < c,
\]
and
\[
f(x) > f(c) \quad \text{for all } x \in I \text{ with } c < x < c + \delta.
\]
\hyperref[prf:positive-derivative-implies-locally-increasing]{Go to proof.}
\end{lemma}
\end{lemmabox}

\begin{remark*}[Standard quantified statement]
\[
\operatorname{Derivative}(L,f,c,\mathbb{R},\mathbb{R}) \;\land\; L > 0
\;\Rightarrow\;
\exists \delta > 0\;\forall x \in I\;
\bigl(c-\delta < x < c \Rightarrow f(x) < f(c)\bigr)
\;\land\;
\bigl(c < x < c+\delta \Rightarrow f(x) > f(c)\bigr).
\]
\end{remark*}

\begin{remark*}[Interpretation]
\[
\forall \delta > 0\;\exists x \in I\;
\bigl((c-\delta < x < c \;\land\; f(x) \geq f(c))
\;\lor\;
(c < x < c+\delta \;\land\; f(x) \leq f(c))\bigr)
\;\Rightarrow\;
f'(c) \leq 0.
\]


\[
\forall \delta > 0\;\exists x \in I\;
\bigl((c-\delta < x < c \;\land\; f(x) \geq f(c))
\;\lor\;
(c < x < c+\delta \;\land\; f(x) \leq f(c))\bigr)
\;\Rightarrow\;
f'(c) \leq 0.
\]


If the instantaneous rate of change is strictly positive at $c$, the function is genuinely climbing through $c$: it is strictly below $f(c)$ immediately to the left and strictly above immediately to the right. The proof is a direct $\varepsilon$-$\delta$ argument: take $\varepsilon = f'(c)/2 > 0$; the definition of the derivative supplies $\delta$ such that the difference quotient is within $\varepsilon$ of $f'(c)$, forcing it to be positive, which then forces the correct sign on $f(x) - f(c)$. This lemma is the key ingredient in the proof of the First Derivative Test.
\end{remark*}

\begin{dependencies}
\begin{itemize}
  \item \hyperref[def:derivative-at-a-point]{Definition: Derivative at a Point}
\end{itemize}
\end{dependencies}

% ---------------------------------------------------------

\begin{lemmabox}{Lemma}
\begin{lemma}[Negative Derivative Implies Locally Decreasing]
\label{lem:negative-derivative-implies-locally-decreasing}
Let $I \subseteq \mathbb{R}$ be an interval, let $f : I \to \mathbb{R}$, let $c \in I$, and suppose that $f$ is differentiable at $c$. If $f'(c) < 0$, then there exists $\delta > 0$ such that
\[
f(x) > f(c) \quad \text{for all } x \in I \text{ with } c - \delta < x < c,
\]
and
\[
f(x) < f(c) \quad \text{for all } x \in I \text{ with } c < x < c + \delta.
\]
\hyperref[prf:negative-derivative-implies-locally-decreasing]{Go to proof.}
\end{lemma}
\end{lemmabox}

\begin{remark*}[Standard quantified statement]
\[
\operatorname{Derivative}(L,f,c,\mathbb{R},\mathbb{R}) \;\land\; L < 0
\;\Rightarrow\;
\exists \delta > 0\;\forall x \in I\;
\bigl(c-\delta < x < c \Rightarrow f(x) > f(c)\bigr)
\;\land\;
\bigl(c < x < c+\delta \Rightarrow f(x) < f(c)\bigr).
\]
\end{remark*}

\begin{remark*}[Interpretation]
\[
\forall \delta > 0\;\exists x \in I\;
\bigl((c-\delta < x < c \;\land\; f(x) \leq f(c))
\;\lor\;
(c < x < c+\delta \;\land\; f(x) \geq f(c))\bigr)
\;\Rightarrow\;
f'(c) \geq 0.
\]


\[
\neg\exists \delta > 0\;\forall x \in I\;
\bigl(c-\delta < x < c \Rightarrow f(x) > f(c)\bigr)
\;\land\;
\bigl(c < x < c+\delta \Rightarrow f(x) < f(c)\bigr)
\;\Rightarrow\;
f'(c) \geq 0.
\]


If the derivative at $c$ is strictly negative, nearby values move downward as the input passes through $c$ from left to right.
\end{remark*}

\begin{dependencies}
\begin{itemize}
  \item \hyperref[lem:positive-derivative-implies-locally-increasing]{Lemma: Positive Derivative Implies Locally Increasing}
\end{itemize}
\end{dependencies}

% ---------------------------------------------------------

\begin{theorembox}{Theorem (First Derivative Test: Relative Maximum)}
\begin{theorem}[First Derivative Test: Relative Maximum]
\label{thm:first-derivative-test-maximum}
Let $I := [a,b] \subseteq \mathbb{R}$, let $f : I \to \mathbb{R}$ be continuous on $I$, and let $c \in (a,b)$. Suppose $f$ is differentiable on $(a,c)$ and $(c,b)$. If there exists $\delta > 0$ such that $(c-\delta,c+\delta) \subseteq I$ and
\[
f'(x) \geq 0 \quad \text{for all } x \in (c-\delta,c),
\qquad
f'(x) \leq 0 \quad \text{for all } x \in (c,c+\delta),
\]
then $f$ has a relative maximum at $c$.
\hyperref[prf:first-derivative-test-maximum]{Go to proof.}
\end{theorem}
\end{theorembox}

\begin{remark*}[Standard quantified statement]
\begin{align*}
&\exists \delta > 0\;\bigl[(c-\delta,c+\delta)\subseteq I
\;\land\;
\forall x\in(c-\delta,c)\;(f'(x)\geq 0)
\;\land\;
\forall x\in(c,c+\delta)\;(f'(x)\leq 0)\bigr]\\
&\quad\Rightarrow\; \operatorname{IsLocalMaximum}(f,c,A).
\end{align*}
\end{remark*}

\begin{remark*}[Predicate reading]
If $f'$ is nonnegative just to the left of $c$ and nonpositive just to the
right of $c$, then $c$ is a local maximum of $f$.
\[
  f'\colon A\to\mathbb{R},\\

  f\colon A\to\mathbb{R},\\
  \operatorname{Derivative}(D,f,a,\mathbb{R},\mathbb{R}).
\]
\end{remark*}

\begin{remark*}[Interpretation]
If the function is nondecreasing immediately to the left of $c$ and nonincreasing immediately to the right, then $c$ is a local maximum --- the function climbs to $c$ and then falls away. The proof integrates: by the nondecreasing characterisation applied on $(c-\delta,c)$, $f(x) \leq f(c)$ for $x < c$; by the nonincreasing characterisation on $(c,c+\delta)$, $f(x) \leq f(c)$ for $x > c$. The sign condition on $f'$ replaces the Interior Extremum Theorem's converse, which does not hold in general; this is a sufficient condition, not a necessary one. The test is applicable even when $f'(c) = 0$ or $f'(c)$ does not exist.
\end{remark*}

\begin{dependencies}
\begin{itemize}
  \item \hyperref[def:relative-maximum]{Definition: Relative Maximum}
  \item \hyperref[thm:nondecreasing-iff-nonneg-derivative]{Theorem: Nondecreasing Iff Nonnegative Derivative}
  \item \hyperref[thm:nonincreasing-iff-nonpos-derivative]{Theorem: Nonincreasing Iff Nonpositive Derivative}
\end{itemize}
\end{dependencies}

% ---------------------------------------------------------

\begin{theorembox}{Theorem}
\begin{theorem}[First Derivative Test: Relative Minimum]
\label{thm:first-derivative-test-minimum}
Let $I := [a,b] \subseteq \mathbb{R}$, let $f : I \to \mathbb{R}$ be continuous on $I$, and let $c \in (a,b)$. Suppose $f$ is differentiable on $(a,c)$ and $(c,b)$. If there exists $\delta > 0$ such that $(c-\delta,c+\delta) \subseteq I$ and
\[
f'(x) \leq 0 \quad \text{for all } x \in (c-\delta,c),
\qquad
f'(x) \geq 0 \quad \text{for all } x \in (c,c+\delta),
\]
then $f$ has a relative minimum at $c$.
\hyperref[prf:first-derivative-test-minimum]{Go to proof.}
\end{theorem}
\end{theorembox}

\begin{remark*}[Standard quantified statement]
\begin{align*}
&\exists \delta > 0\;\bigl[(c-\delta,c+\delta)\subseteq I
\;\land\;
\forall x\in(c-\delta,c)\;(f'(x)\leq 0)
\;\land\;
\forall x\in(c,c+\delta)\;(f'(x)\geq 0)\bigr]\\
&\quad\Rightarrow\; \operatorname{IsLocalMinimum}(f,c,A).
\end{align*}
\end{remark*}

\begin{remark*}[Predicate reading]
If $f'$ is nonpositive just to the left of $c$ and nonnegative just to the
right of $c$, then $c$ is a local minimum of $f$.
\[
  f'\colon A\to\mathbb{R},\\

  f\colon A\to\mathbb{R},\\
  \operatorname{Derivative}(D,f,a,\mathbb{R},\mathbb{R}).
\]
\end{remark*}

\begin{remark*}[Interpretation]
If the derivative changes from nonpositive to nonnegative at $c$, the graph falls into $c$ and then rises away from it, so $c$ is a local minimum.
\end{remark*}

\begin{dependencies}
\begin{itemize}
  \item \hyperref[def:relative-minimum]{Definition: Relative Minimum}
  \item \hyperref[thm:first-derivative-test-maximum]{Theorem: First Derivative Test: Relative Maximum}
\end{itemize}
\end{dependencies}


\subsection{Second-Order Shape}

\begin{toolkitbox}{Second-Order Shape --- Quick Reference}
\small
\begin{tabular}{l l}
\toprule
\textbf{Concept} & \textbf{Role} \\
\midrule
\hyperref[def:second-derivative]{Second derivative} & Derivative of the derivative \\
\hyperref[def:higher-derivatives]{Higher derivatives} & Iterated derivatives \(f^{(n)}\) \\
\hyperref[thm:second-derivative-convexity-test]{Convexity test} & \(f''\geq 0\) detects convexity \\
\hyperref[thm:second-derivative-concavity-test]{Concavity test} & \(f''\leq 0\) detects concavity \\
\hyperref[thm:second-derivative-test]{Second derivative test} & Classifies critical points \\
\hyperref[prop:inflection-point-necessary-condition]{Inflection condition} & Continuous \(f''\) changes sign through zero \\
\bottomrule
\end{tabular}
\end{toolkitbox}

\subsection{Higher-Order Derivatives}

\begin{definitionbox}{Definition (Second Derivative)}
\begin{definition}[Second Derivative]
\label{def:second-derivative}
Let $f$ be differentiable on a neighbourhood of $c$. If $f'$ is differentiable at $c$, the \emph{second derivative} of $f$ at $c$ is
\[
  f''(c):=(f')'(c).
\]
\end{definition}
\end{definitionbox}

\begin{remark*}[Standard quantified statement]
\[
  \operatorname{IsDifferentiable}(f',c,f''(c))
  \quad\Longleftrightarrow\quad
  f''(c)=(f')'(c).
\]
\end{remark*}

\begin{remark*}[Interpretation]
The first derivative measures instantaneous slope. The second derivative measures how that slope itself changes.
\end{remark*}

\begin{dependencies}
\begin{itemize}
  \item \hyperref[def:derivative-at-a-point]{Derivative at a Point}
\end{itemize}
\end{dependencies}

\begin{definitionbox}{Definition (Higher Derivatives)}
\begin{definition}[Higher Derivatives]
\label{def:higher-derivatives}
Define $f^{(0)}:=f$ and $f^{(1)}:=f'$. For $n\geq 2$, if $f^{(n-1)}$ is differentiable at $c$, define
\[
  f^{(n)}(c):=(f^{(n-1)})'(c).
\]
\end{definition}
\end{definitionbox}

\begin{remark*}[Standard quantified statement]
\[
  \forall n\geq 2:\quad
  f^{(n)}(c)=(f^{(n-1)})'(c)
\]
whenever the derivative on the right exists.
\end{remark*}

\begin{remark*}[Interpretation]
Higher derivatives are obtained by iterating the derivative operator. Their existence is not automatic: each stage requires differentiability of the previous derivative.
\end{remark*}

\begin{dependencies}
\begin{itemize}
  \item \hyperref[def:second-derivative]{Second Derivative}
\end{itemize}
\end{dependencies}

\subsection{Second-Order Shape}

\begin{theorembox}{Theorem (Second-Derivative Convexity Test)}
\begin{theorem}[Second-Derivative Convexity Test]
\label{thm:second-derivative-convexity-test}
Let $I\subseteq\mathbb{R}$ be an interval and let $f:I\to\mathbb{R}$ be twice differentiable on $I$. If
\[
  f''(x)\geq 0 \qquad \text{for all } x\in I,
\]
then $f$ is convex on $I$.
\hyperref[prf:second-derivative-convexity-test]{Go to proof.}
\end{theorem}
\end{theorembox}

\begin{remark*}[Standard quantified statement]
\[
  \bigl(\forall x\in I,\ f''(x)\geq 0\bigr)
  \Longrightarrow
  \operatorname{IsConvex}(f,I).
\]
\end{remark*}

\begin{remark*}[Interpretation]
Nonnegative second derivative means the slopes do not decrease. This is the derivative-based signature of convex bending.
\end{remark*}

\begin{dependencies}
\begin{itemize}
  \item \hyperref[def:convex-function]{Convex Function}
  \item \hyperref[def:second-derivative]{Second Derivative}
\end{itemize}
\end{dependencies}

\begin{theorembox}{Theorem (Second-Derivative Concavity Test)}
\begin{theorem}[Second-Derivative Concavity Test]
\label{thm:second-derivative-concavity-test}
Let $I\subseteq\mathbb{R}$ be an interval and let $f:I\to\mathbb{R}$ be twice differentiable on $I$. If
\[
  f''(x)\leq 0 \qquad \text{for all } x\in I,
\]
then $f$ is concave on $I$.
\hyperref[prf:second-derivative-concavity-test]{Go to proof.}
\end{theorem}
\end{theorembox}

\begin{remark*}[Standard quantified statement]
\[
  \bigl(\forall x\in I,\ f''(x)\leq 0\bigr)
  \Longrightarrow
  \operatorname{IsConcave}(f,I).
\]
\end{remark*}

\begin{remark*}[Interpretation]
Nonpositive second derivative means the slopes do not increase. This is the derivative-based signature of concave bending.


The second-derivative tests are narrower than the chord definitions of convexity and concavity. The function $|x|$ is convex, but it is not twice differentiable at $0$.
\end{remark*}

\begin{dependencies}
\begin{itemize}
  \item \hyperref[def:concave-function]{Concave Function}
  \item \hyperref[def:second-derivative]{Second Derivative}
\end{itemize}
\end{dependencies}

\begin{figure}[H]
\centering
\resizebox{0.82\linewidth}{!}{%
\begin{tikzpicture}[scale=0.82]
  \def\mn{0.7*(\x-1.1)*(\x-1.1)+0.5}
  \glowcurve{color=atlasteal, domain=0.1:2.1, axis=0, top=3.0, expr={\mn}}
  \node[atlas dot] at (1.1,0.5){};
  \draw[dropline] (1.1,0.5)--(1.1,0);
  \node[atlas label,below] at (1.1,-0.05){\footnotesize$c$};
  \node[atlas label,below,align=center] at (1.0,-0.40)
    {\tiny$f'(c)=0$ and\\[-1pt]$f''(c)>0\Rightarrow$ min};
  \begin{scope}[xshift=3.4cm]
    \def\mx{2.6-0.7*(\x-1.1)*(\x-1.1)}
    \glowcurve{color=atlasorange, domain=0.1:2.1, axis=0, top=3.0, expr={\mx}}
    \node[atlas dot] at (1.1,2.6){};
    \draw[dropline] (1.1,2.6)--(1.1,0);
    \node[atlas label,below] at (1.1,-0.05){\footnotesize$c$};
    \node[atlas label,below,align=center] at (1.2,-0.40)
      {\tiny$f'(c)=0$ and\\[-1pt]$f''(c)<0\Rightarrow$ max};
  \end{scope}
\end{tikzpicture}%
}
\caption{Second derivative test: the sign of \(f''(c)\) decides the local shape at a critical point.}
\label{fig:second-derivative-tests}
\end{figure}


\begin{theorembox}{Theorem (Second Derivative Test)}
\begin{theorem}[Second Derivative Test]
\label{thm:second-derivative-test}
Let $f$ be twice differentiable in a neighbourhood of $c$, and suppose $f'(c)=0$.
\begin{itemize}
  \item If $f''(c)>0$, then $f$ has a strict relative minimum at $c$.
  \item If $f''(c)<0$, then $f$ has a strict relative maximum at $c$.
\end{itemize}
\hyperref[prf:second-derivative-test]{Go to proof.}
\end{theorem}
\end{theorembox}

\begin{remark*}[Standard quantified statement]
\[
  f'(c)=0\land f''(c)>0 \Rightarrow \operatorname{IsStrictRelativeMinimum}(f,c),
\]
and
\[
  f'(c)=0\land f''(c)<0 \Rightarrow \operatorname{IsStrictRelativeMaximum}(f,c).
\]
\end{remark*}

\begin{remark*}[Interpretation]
At a critical point, positive second derivative means the graph bends upward; negative second derivative means it bends downward.
\end{remark*}

\begin{dependencies}
\begin{itemize}
  \item \hyperref[def:second-derivative]{Second Derivative}
  \item \hyperref[def:relative-minimum]{Relative Minimum}
  \item \hyperref[def:relative-maximum]{Relative Maximum}
\end{itemize}
\end{dependencies}

\begin{propositionbox}{Proposition (Inflection Point Necessary Condition)}
\begin{proposition}[Inflection Point Necessary Condition]
\label{prop:inflection-point-necessary-condition}
Let $f$ have an inflection point at $c$. If $f''$ exists on a neighbourhood of $c$ and is continuous at $c$, then
\[
  f''(c)=0.
\]
\hyperref[prf:inflection-point-necessary-condition]{Go to proof.}
\end{proposition}
\end{propositionbox}

\begin{remark*}[Standard quantified statement]
\[
  \operatorname{IsInflection}(f,c)\land \operatorname{IsContinuous}(f'',c)
  \Longrightarrow
  f''(c)=0.
\]
\end{remark*}

\begin{remark*}[Interpretation]
An inflection point is a change in bending. When the second derivative is continuous, such a change can pass through $c$ only by passing through zero. The converse can fail: $f''(c)=0$ does not by itself guarantee an inflection point.
\end{remark*}

\begin{dependencies}
\begin{itemize}
  \item \hyperref[def:inflection-point]{Inflection Point}
  \item \hyperref[def:second-derivative]{Second Derivative}
  \item \hyperref[def:continuous-at-point]{Continuity at a Point}
\end{itemize}
\end{dependencies}


\subsection{Darboux's Theorem}

\begin{toolkitbox}{Darboux --- Quick Reference}
\small
\begin{tabular}{l l}
\toprule
\textbf{Statement} & \textbf{Role} \\
\midrule
\hyperref[thm:darboux]{Darboux's Theorem} & Derivatives have the intermediate value property \\
\bottomrule
\end{tabular}
\end{toolkitbox}

% =========================================================
% Section: Darboux's Theorem
% =========================================================
% ---------------------------------------------------------

\begin{theorembox}{Theorem (Darboux's Theorem)}
\begin{theorem}[Darboux's Theorem]
\label{thm:darboux}
Let $I := [a,b] \subseteq \mathbb{R}$ and let $f : I \to \mathbb{R}$ be differentiable on $I$. If $k$ is a real number strictly between $f'(a)$ and $f'(b)$, then there exists $c \in (a,b)$ such that
\[
f'(c) = k.
\]
\hyperref[prf:darboux]{Go to proof.}
\end{theorem}
\end{theorembox}

\begin{remark*}[Standard quantified statement]
\[
\forall k\;\bigl(\min(f'(a),f'(b)) < k < \max(f'(a),f'(b))
\;\Rightarrow\; \exists c \in (a,b)\;(f'(c)=k)\bigr).
\]
\end{remark*}

\begin{remark*}[Predicate reading]
The displayed quantified statement is the predicate-level form of $label: its hypotheses name the input conditions, and its conclusion names the asserted structure or relation.
\[
  f\colon A\to\mathbb{R},\\
  \operatorname{Derivative}(D,f,a,\mathbb{R},\mathbb{R}).
\]
\end{remark*}

\begin{remark*}[Interpretation]
The Intermediate Value Theorem says continuous functions have the intermediate
value property. Darboux's theorem says derivatives also have that property even
when the derivative is not continuous. Thus a derivative cannot have a jump
discontinuity. The proof uses differentiability of $f$ to construct an
auxiliary function $g(x)=f(x)-kx$ and then applies the Extreme Value Theorem
and the Interior Extremum Theorem.
\end{remark*}

\begin{dependencies}
\begin{itemize}
  \item \hyperref[def:derivative-at-a-point]{Definition: Derivative at a Point}
  \item \hyperref[thm:extreme-value-theorem]{Theorem: Extreme Value Theorem}
  \item \hyperref[thm:interior-extremum-theorem]{Theorem: Interior Extremum Theorem}
\end{itemize}
\end{dependencies}


\subsection{Smoothness Classes}

\begin{toolkitbox}{Smoothness Classes --- Quick Reference}
\small
\begin{tabular}{l l}
\toprule
\textbf{Concept} & \textbf{Meaning} \\
\midrule
\hyperref[def:class-c1]{Class \(C^1\)} & Differentiable with continuous first derivative \\
\hyperref[def:class-ck]{Class \(C^k\)} & \(k\) derivatives with continuous top derivative \\
\hyperref[def:class-cinfty]{Class \(C^\infty\)} & Continuous derivatives of every order \\
\hyperref[def:class-comega]{Class \(C^\omega\)} & Equal to local Taylor series \\
\hyperref[thm:smoothness-tower]{Smoothness tower} & Strict hierarchy \(C^\omega\subsetneq C^\infty\subsetneq\cdots\) \\
\hyperref[def:class-c11]{Class \(C^{1,1}\)} & \(C^1\) with Lipschitz first derivative \\
\hyperref[thm:c11-placement]{Placement of \(C^{1,1}\)} & Refines the gap between \(C^1\) and \(C^2\) \\
\hyperref[cor:bounded-second-derivative-implies-c11]{Bounded \(f''\)} & Bounded second derivative gives \(C^{1,1}\) \\
\hyperref[ex:x2sin]{Example: differentiable not \(C^1\)} & Oscillatory derivative discontinuity at \(0\) \\
\bottomrule
\end{tabular}
\end{toolkitbox}

\begin{definitionbox}{Definition (Continuously Differentiable; Class \(C^1\))}
\begin{definition}[Continuously Differentiable; Class \(C^1\)]
\label{def:class-c1}
Let \(I\subseteq\mathbb{R}\) be an interval and let \(f:I\to\mathbb{R}\). The function \(f\) is \emph{continuously differentiable} on \(I\), written \(f\in C^1(I)\), if \(f\) is differentiable at every point of \(I\) and the derivative \(f':I\to\mathbb{R}\) is continuous on \(I\).
\end{definition}
\end{definitionbox}

\begin{remark*}[Standard quantified statement]
\[
  f\in C^1(I)
  \Longleftrightarrow
  \bigl(\forall x\in I:\operatorname{IsDifferentiable}(f,x)\bigr)
  \land
  \operatorname{IsContinuous}(f',I).
\]
\end{remark*}

\begin{remark*}[Interpretation]
The condition \(f\in C^1(I)\) is strictly stronger than differentiability. A function may be differentiable everywhere while its derivative fails to be continuous; Darboux's Theorem still constrains that failure, because derivatives cannot jump. Thus the gap between differentiability and \(C^1\) is oscillatory rather than jump-like.
\end{remark*}

\begin{dependencies}
\begin{itemize}
  \item \hyperref[def:derivative-at-a-point]{Definition: Derivative at a Point}
  \item \hyperref[def:continuous-at-point]{Definition: Continuity at a Point}
\end{itemize}
\end{dependencies}

\begin{exposition}
\phantomsection\label{ex:x2sin}
Let
\[
  f(x)=
  \begin{cases}
    x^2\sin(1/x), & x\neq 0,\\
    0, & x=0.
  \end{cases}
\]
At the origin,
\[
  f'(0)=\lim_{x\to0}\frac{x^2\sin(1/x)}{x}
  =
  \lim_{x\to0}x\sin(1/x)=0
\]
by squeezing. For \(x\neq0\),
\[
  f'(x)=2x\sin(1/x)-\cos(1/x).
\]
The first term tends to \(0\), but \(\cos(1/x)\) oscillates, so \(f'\) is
not continuous at \(0\). Thus \(f\) is differentiable on all of
\(\mathbb{R}\), but \(f\notin C^1(\mathbb{R})\). The failure is oscillatory,
not a jump, as Darboux's theorem requires.
\end{exposition}

\begin{definitionbox}{Definition (Class \(C^k\))}
\begin{definition}[Class \(C^k\)]
\label{def:class-ck}
Let \(I\subseteq\mathbb{R}\) be an interval, let \(k\in\mathbb{N}\), and let \(f:I\to\mathbb{R}\). The function \(f\) is of \emph{class \(C^k\)} on \(I\), written \(f\in C^k(I)\), if the derivatives \(f^{(0)},f^{(1)},\ldots,f^{(k)}\) exist on \(I\) and \(f^{(k)}\) is continuous on \(I\). By convention, \(C^0(I)\) is the class of continuous functions on \(I\).
\end{definition}
\end{definitionbox}

\begin{remark*}[Standard quantified statement]
\[
  f\in C^k(I)
  \Longleftrightarrow
  \bigl(\forall j\in\{0,1,\ldots,k\}: f^{(j)}\text{ exists on }I\bigr)
  \land
  \operatorname{IsContinuous}\bigl(f^{(k)},I\bigr).
\]
\end{remark*}

\begin{remark*}[Predicate reading]
\[
  \operatorname{IsClassC}(f,I,k).
\]
For a fixed interval \(I\), the classes are nested:
\[
  f\in C^{k+1}(I)\Longrightarrow f\in C^k(I).
\]
\end{remark*}

\begin{remark*}[Negated quantified statement]
\[
  f\notin C^k(I)
  \Longleftrightarrow
  \exists x\in I:\ \bigl[f^{(k)}\text{ fails to exist at }x\bigr]
  \vee
  \bigl[f^{(k)}\text{ exists on }I\land
  \neg\operatorname{IsContinuous}(f^{(k)},x)\bigr].
\]
\end{remark*}

\begin{remark*}[Negation predicate reading]
\(\neg\operatorname{IsClassC}(f,I,k)\): somewhere the \(k\)-th derivative
either does not exist or exists but is not continuous. A single point where
the top derivative is discontinuous certifies \(f\notin C^k(I)\).
\end{remark*}

\begin{remark*}[Interpretation]
The notation \(C^k(I)\) records \(k\) layers of differentiability together with continuity at the top layer. This is the regularity scale used to state how many times a theorem may differentiate a function without leaving the controlled setting.
\end{remark*}

\begin{dependencies}
\begin{itemize}
  \item \hyperref[def:higher-derivatives]{Definition: Higher Derivatives}
  \item \hyperref[def:class-c1]{Definition: Continuously Differentiable; Class \(C^1\)}
\end{itemize}
\end{dependencies}

\begin{definitionbox}{Definition (Smooth Function; Class \(C^\infty\))}
\begin{definition}[Smooth Function; Class \(C^\infty\)]
\label{def:class-cinfty}
Let \(I\subseteq\mathbb{R}\) be an interval and let \(f:I\to\mathbb{R}\). The function \(f\) is \emph{smooth}, or of \emph{class \(C^\infty\)}, written \(f\in C^\infty(I)\), if \(f\in C^k(I)\) for every \(k\in\mathbb{N}\). Equivalently,
\[
  C^\infty(I)=\bigcap_{k=0}^{\infty}C^k(I).
\]
\end{definition}
\end{definitionbox}

\begin{remark*}[Standard quantified statement]
\[
  f\in C^\infty(I)
  \Longleftrightarrow
  \forall k\in\mathbb{N}:\ f\in C^k(I).
\]
\end{remark*}

\begin{remark*}[Negated quantified statement]
\[
  f\notin C^\infty(I)
  \Longleftrightarrow
  \exists k\in\mathbb{N}:\ f\notin C^k(I).
\]
A single order at which differentiability or continuity breaks demotes the
function out of \(C^\infty\).
\end{remark*}

\begin{remark*}[Interpretation]
Smoothness is the regularity language imported by differential geometry. Smooth manifolds are assembled from \(C^\infty\) transition maps, and the tangent-map construction seeded by the differential uses this vocabulary.
\end{remark*}

\begin{dependencies}
\begin{itemize}
  \item \hyperref[def:class-ck]{Definition: Class \(C^k\)}
\end{itemize}
\end{dependencies}

\begin{definitionbox}{Definition (Real-Analytic Function; Class \(C^\omega\))}
\begin{definition}[Real-Analytic Function; Class \(C^\omega\)]
\label{def:class-comega}
Let \(I\subseteq\mathbb{R}\) be an interval and let \(f:I\to\mathbb{R}\). The function \(f\) is \emph{real-analytic}, or of \emph{class \(C^\omega\)}, written \(f\in C^\omega(I)\), if for every \(a\in I\) there exists \(r>0\) such that
\[
  f(x)=\sum_{n=0}^{\infty}\frac{f^{(n)}(a)}{n!}(x-a)^n
  \qquad\text{for every }x\in(a-r,a+r)\cap I.
\]
\end{definition}
\end{definitionbox}

\begin{remark*}[Standard quantified statement]
\[
  f\in C^\omega(I)
  \Longleftrightarrow
  \forall a\in I\;\exists r>0\;\forall x\in(a-r,a+r)\cap I:\;
  f(x)=\sum_{n=0}^{\infty}\frac{f^{(n)}(a)}{n!}(x-a)^n.
\]
\end{remark*}

\begin{remark*}[Predicate reading]
The displayed quantified statement is the predicate-level form of $label: its hypotheses name the input conditions, and its conclusion names the asserted structure or relation.
\end{remark*}

\begin{remark*}[Interpretation]
Analytic regularity means that the Taylor series does more than approximate the function: near each point, the series equals the function. The infinite Taylor series and its convergence are developed in \hyperref[def:taylor-polynomial-at-a-point]{Taylor's construction and error analysis}, which supplies the series machinery behind this definition.
\end{remark*}

\begin{dependencies}
\begin{itemize}
  \item \hyperref[def:class-cinfty]{Definition: Smooth Function; Class \(C^\infty\)}
  \item \hyperref[def:taylor-polynomial-at-a-point]{Definition: Taylor Polynomial at a Point}
\end{itemize}
\end{dependencies}

\begin{theorembox}{Theorem (The Smoothness Tower)}
\begin{theorem}[The Smoothness Tower]
\label{thm:smoothness-tower}
For every nondegenerate interval \(I\), the inclusions
\[
  C^0(I)\supseteq C^1(I)\supseteq C^2(I)\supseteq \cdots
  \supseteq C^\infty(I)\supseteq C^\omega(I)
\]
hold, and each displayed inclusion is strict.
\hyperref[prf:smoothness-tower]{Go to proof.}
\end{theorem}
\end{theorembox}

\begin{remark*}[Standard quantified statement]
\[
  \forall k\in\mathbb{N}:\ C^{k+1}(I)\subsetneq C^k(I),
  \qquad
  C^\omega(I)\subsetneq C^\infty(I).
\]
\end{remark*}

\begin{remark*}[Failure modes]
\begin{description}
  \item[Exposition.]
The inclusions do not fail; only their strictness needs witnesses.
\begin{itemize}
  \item \(C^k\setminus C^{k+1}\): \(\varphi_k(x)=x^k|x|\), whose \(k\)-th
        derivative is \((k+1)!\,|x|\), continuous but not differentiable at
        \(0\).
  \item \(C^\infty\setminus C^\omega\): the flat function
        \(e^{-1/x^2}\), smooth with all derivatives zero at \(0\), so its
        Maclaurin series is zero while the function is not.
\end{itemize}
These separators prove every displayed containment is a genuine drop.
\end{description}
\end{remark*}

\begin{remark*}[Interpretation]
The inclusions are definitional; strictness is witnessed by examples. On any nondegenerate interval, translated and rescaled copies of \(x\mapsto x^k|x|\) separate \(C^k\) from \(C^{k+1}\). Translated copies of the flat function
\[
  f(x)=
  \begin{cases}
    e^{-1/x^2}, & x\neq 0,\\
    0, & x=0,
  \end{cases}
\]
separate \(C^\infty\) from \(C^\omega\): all derivatives at the flat point vanish, so the Taylor series is identically zero, while the function is positive on one side of that point.


Between \(C^1\) and \(C^2\) there is a regularity floor worth naming. A
\(C^1\) function has a continuous derivative; a \(C^2\) function has a
derivative that is itself differentiable. Many functions used in analysis
live in the gap: their first derivative is continuous and controlled, but a
classical second derivative may fail to exist at a few points. The control is
a Lipschitz bound on \(f'\), and the class is written \(C^{1,1}\).
\end{remark*}

\begin{dependencies}
\begin{itemize}
  \item \hyperref[def:class-ck]{Definition: Class \(C^k\)}
  \item \hyperref[def:class-cinfty]{Definition: Smooth Function; Class \(C^\infty\)}
  \item \hyperref[def:class-comega]{Definition: Real-Analytic Function; Class \(C^\omega\)}
  \item \hyperref[thm:darboux]{Theorem: Darboux's Theorem}
\end{itemize}
\end{dependencies}

\begin{figure}[H]
\centering
\resizebox{0.76\linewidth}{!}{%
\begin{tikzpicture}[scale=0.82]
  \def\lf{2.0+1.05*sin(80*\x)*0.9+0.16*\x}
  \glowcurve{color=atlasblue, domain=0.0:4.2, axis=0, top=4.2, expr={\lf}}
  \draw[atlas axes] (0,0)--(4.8,0);
  \draw[atlas axes] (0,0)--(0,4.4);
  \def\L{0.85}
  \draw[atlas probe,atlasconegreen] (0,3.6)--(4.2,{3.6-\L*4.2});
  \draw[atlas probe,atlasconegreen] (0,0.5)--(4.2,{0.5+\L*4.2});
  \draw[atlas probe,atlasconegreen] (0,1.0)--(4.2,{1.0+\L*4.2});
  \draw[atlas probe,atlasconegreen] (0,3.1)--(4.2,{3.1-\L*4.2});
  \node[atlas label,above] at (2.1,4.15){\scriptsize$|f'(x)-f'(y)|\le L|x-y|$};
  \node[atlas label,left] at (0,3.9){\scriptsize$-L$};
  \node[atlas label,left] at (0,0.5){\scriptsize$\pm L$};
  \node[atlas label,right] at (4.2,4.0){\scriptsize$\pm L$};
  \node[atlas label,right] at (4.2,0.6){\scriptsize$\pm L$};
  \node[atlas label,right] at (3.4,3.1){\scriptsize$y=f(x)$};
  \node[atlas label] at (2.5,1.0){\scriptsize$-L$};
\end{tikzpicture}%
}
\caption{A Lipschitz derivative changes at a uniformly bounded rate.}
\label{fig:lipschitz-derivative}
\end{figure}


\begin{definitionbox}{Definition (Lipschitz-Derivative Class \(C^{1,1}\))}
\begin{definition}[Lipschitz-Derivative Class \(C^{1,1}\)]
\label{def:class-c11}
Let \(I\subseteq\mathbb{R}\) be an interval and let \(f:I\to\mathbb{R}\).
The function \(f\) is of \emph{class \(C^{1,1}\)} on \(I\), written
\(f\in C^{1,1}(I)\), if \(f\in C^1(I)\) and \(f'\) is Lipschitz on \(I\):
there exists \(L\geq 0\) such that
\[
  |f'(x)-f'(y)|\leq L|x-y|
  \qquad\text{for all }x,y\in I.
\]
\end{definition}
\end{definitionbox}

\begin{remark*}[Standard quantified statement]
\[
  f\in C^{1,1}(I)
  \Longleftrightarrow
  f\in C^1(I)\land
  \exists L\geq 0\;\forall x,y\in I:\ |f'(x)-f'(y)|\leq L|x-y|.
\]
\end{remark*}

\begin{remark*}[Predicate reading]
\[
  f\in C^{1,1}(I)
  \Longleftrightarrow
  \operatorname{IsClassC}(f,I,1)\land
  \exists L\geq0:\operatorname{IsLipschitz}(f',I,L).
\]
The smallest admissible \(L\) is the Lipschitz constant of \(f'\):
\[
  \sup_{x\neq y}\frac{|f'(x)-f'(y)|}{|x-y|}.
\]
\end{remark*}

\begin{remark*}[Negated quantified statement]
\[
  f\notin C^{1,1}(I)
  \Longleftrightarrow
  f\notin C^1(I)\vee
  \bigl(\forall L\geq0\;\exists x,y\in I:\ |f'(x)-f'(y)|>L|x-y|\bigr).
\]
\end{remark*}

\begin{remark*}[Interpretation]
\(C^{1,1}\) says that \(f'\) behaves as if \(|f''|\leq L\), without requiring
\(f''\) to exist everywhere. It is the exact regularity tracked by many
first-order estimates: the derivative is not merely continuous, but changes
at a bounded rate.
\end{remark*}

\begin{dependencies}
\begin{itemize}
  \item \hyperref[def:class-c1]{Definition: Continuously Differentiable; Class \(C^1\)}
  \item \hyperref[cor:derivative-bound-implies-lipschitz]{Corollary: Derivative Bound Implies Lipschitz}
\end{itemize}
\end{dependencies}

\begin{theorembox}{Theorem (Placement of \(C^{1,1}\))}
\begin{theorem}[Placement of \(C^{1,1}\)]
\label{thm:c11-placement}
For every nondegenerate interval \(I\),
\[
  C^{1,1}(I)\subsetneq C^1(I),
\]
and \(C^{1,1}(I)\) is not contained in \(C^2(I)\). If \(K\) is a compact
interval, then
\[
  C^2(K)\subseteq C^{1,1}(K)\subsetneq C^1(K),
\]
and the first inclusion is strict.
\hyperref[prf:c11-placement]{Go to proof.}
\end{theorem}
\end{theorembox}

\begin{remark*}[Standard quantified statement]
\[
  \bigl(f\in C^{1,1}(I)\Rightarrow f\in C^1(I)\bigr),
  \qquad
  \bigl(K\text{ compact}\land f\in C^2(K)\Rightarrow f\in C^{1,1}(K)\bigr),
\]
with the displayed containments strict.
\end{remark*}

\begin{remark*}[Failure modes]
\begin{description}
  \item[Exposition.]
The global inclusion \(C^2(I)\subseteq C^{1,1}(I)\) is false on arbitrary
intervals: \(f(x)=x^3\) is \(C^2(\mathbb{R})\), but
\(f'(x)=3x^2\) is not Lipschitz on \(\mathbb{R}\). Strictness is witnessed by
two elementary functions.
\begin{itemize}
  \item \(C^{1,1}\setminus C^2\): \(f(x)=x|x|\) has \(f'(x)=2|x|\), which is
        Lipschitz, but \(f''(0)\) does not exist.
  \item \(C^1\setminus C^{1,1}\): \(f(x)=|x|^{4/3}\) has continuous first
        derivative, but \(f'\) has unbounded difference quotients at \(0\).
\end{itemize}
\end{description}
\end{remark*}

\begin{remark*}[Interpretation]
\(C^{1,1}\) is a genuine intermediate regularity condition, but it is not
the same as \(C^2\). A bounded second derivative forces the class; a
Lipschitz first derivative may have corners; and a continuous first
derivative may still vary too sharply to be Lipschitz.
\end{remark*}

\begin{dependencies}
\begin{itemize}
  \item \hyperref[def:class-c11]{Definition: Lipschitz-Derivative Class \(C^{1,1}\)}
  \item \hyperref[thm:smoothness-tower]{Theorem: The Smoothness Tower}
  \item \hyperref[cor:derivative-bound-implies-lipschitz]{Corollary: Derivative Bound Implies Lipschitz}
\end{itemize}
\end{dependencies}

\begin{figure}[H]
\centering
\resizebox{0.78\linewidth}{!}{%
\begin{tikzpicture}[atlas,scale=1.05]
  \def\f{1.8+0.45*sin(70*\x)}
  \glowcurve{color=atlasblue, domain=0.3:4.7, axis=0, top=3.4, expr={\f}}
  \draw[atlas axes] (0,0)--(5.1,0);
  \draw[atlas axes] (0,0)--(0,3.4);
  \def\p{2.0}
  \pgfmathsetmacro\fp{1.8+0.45*sin(70*\p)}
  \pgfmathsetmacro\mp{0.45*70*0.0174533*cos(70*\p)}
  \def\L{0.7}
  \draw[atlas probe,atlasconegreen,dash pattern=on 3pt off 2pt]
    plot[domain=0.45:3.55,samples=80,smooth]
    (\x,{\fp+\mp*(\x-\p)+0.5*\L*(\x-\p)*(\x-\p)});
  \draw[atlas probe,atlasconegreen,dash pattern=on 3pt off 2pt]
    plot[domain=0.45:3.55,samples=80,smooth]
    (\x,{\fp+\mp*(\x-\p)-0.5*\L*(\x-\p)*(\x-\p)});
  \draw[atlas probe,atlasred]
    ({\p-1.4},{\fp-\mp*1.4})--({\p+1.4},{\fp+\mp*1.4});
  \node[atlas dot] at (\p,\fp){};
  \draw[dropline] (\p,\fp)--(\p,0);
  \node[atlas label,below] at (\p,-0.05){\footnotesize$p$};
  \node[atlas label,above] at (2.6,3.15){\scriptsize$|f'(x)-f'(y)|\le L\,|x-y|$};
  \node[atlas label,atlasconegreen!60!atlasink,above left] at (1.0,2.55)
    {\scriptsize$f(p)+f'(p)(x-p)+\tfrac{L}{2}(x-p)^2$};
  \node[atlas label,atlasconegreen!60!atlasink,below right=2pt] at (3.0,0.55)
    {\scriptsize$-\tfrac{L}{2}(x-p)^2$};
  \node[atlas label,atlasred!70!atlasink,below left] at (\p-0.5,{\fp-\mp*0.5})
    {\scriptsize tangent at $p$};
  \node[atlas label,right] at (3.9,2.35){\footnotesize$y=f(x)$};
\end{tikzpicture}%
}
\caption{\(C^{1,1}\) control gives quadratic upper and lower bounds around a tangent line.}
\label{fig:c11-descent-bounds}
\end{figure}


\begin{corollarybox}{Corollary (Bounded Second Derivative Implies \(C^{1,1}\))}
\begin{corollary}[Bounded Second Derivative Implies \(C^{1,1}\)]
\label{cor:bounded-second-derivative-implies-c11}
Let \(I\subseteq\mathbb{R}\) be an interval and let \(f\in C^2(I)\). If
\(|f''(x)|\leq M\) for all \(x\in I\), then \(f\in C^{1,1}(I)\), and \(f'\)
is Lipschitz with constant \(M\).
\hyperref[prf:bounded-second-derivative-implies-c11]{Go to proof.}
\end{corollary}
\end{corollarybox}

\begin{remark*}[Standard quantified statement]
\[
  \bigl(f\in C^2(I)\land \forall x\in I:\ |f''(x)|\leq M\bigr)
  \Rightarrow
  f\in C^{1,1}(I)\land \operatorname{IsLipschitz}(f',I,M).
\]
\end{remark*}

\begin{remark*}[Interpretation]
This is \hyperref[cor:derivative-bound-implies-lipschitz]{Derivative Bound
Implies Lipschitz} applied one level up, to \(f'\) in place of \(f\). The
converse requires almost-everywhere language and is therefore deferred.


A Lipschitz derivative is differentiable at most points, but making that
sentence precise requires measure zero. In one variable this passes through
bounded variation; in several variables it becomes Rademacher's theorem.
Those results belong to the integration and measure-theoretic development,
not to this chapter.
\end{remark*}

\begin{dependencies}
\begin{itemize}
  \item \hyperref[cor:derivative-bound-implies-lipschitz]{Corollary: Derivative Bound Implies Lipschitz}
  \item \hyperref[def:class-ck]{Definition: Class \(C^k\)}
  \item \hyperref[def:class-c11]{Definition: Lipschitz-Derivative Class \(C^{1,1}\)}
\end{itemize}
\end{dependencies}

